\section{Flat dual numbers over ramified discrete valuation rings}
\label{sec:ramified-dual}

Ramification does not restore generation in residue characteristic two.
The cohomology of a suitable integral involution has length one, even
when $O/2O$ has even length.

\begin{theorem}\label{thm:ramified-dual}
Let $O$ be a characteristic-zero discrete valuation ring with residue
characteristic two, and set $A=O[\epsilon]/(\epsilon^2)$.
The kernel of reduction
\[
 K_0(\Lat_A(C_2))\longrightarrow K_0(\Lat_O(C_2))
\]
contains a direct summand isomorphic to a countable direct sum of copies
of $\Z/2$.  In particular $\Phi_A$ is not surjective.  This applies to
$O=\Z_2[\sqrt2]$ and to every ramification index and residue field.
\end{theorem}

\begin{proof}
Choose a uniformizer $\pi$ and put $h=2/\pi\in O$.
For an $A$-free involution write $T=S+\epsilon U$ with matrices over
$O$.  Then $S^2=I$ and $SU+US=0$.  On the two Tate groups
\[
 H^0=\ker(S-I)/\operatorname{im}(S+I),\qquad
 H^1=\ker(S+I)/\operatorname{im}(S-I),
\]
$U$ induces an odd operator $\delta$, interchanging the parities.
These groups have finite $O$-length and are killed by two: the usual
norm and difference maps give the latter assertion, and after
inverting two their complex is exact.

The operator is natural for maps that depend on $\epsilon$.
Indeed, an equivariant map $F=F_0+\epsilon F_1$ satisfies
\[
 S'F_0=F_0S,\qquad U'F_0-F_0U=F_1S-S'F_1.
\]
On a cycle $Sv=\pm v$, the second right-hand side is an appropriate
Tate boundary.  Therefore it induces zero in the next parity.
In an $A$-split exact sequence choose bases with
\[
 S=\begin{pmatrix}S'&a\\0&S''\end{pmatrix},\qquad
 U=\begin{pmatrix}U'&b\\0&U''\end{pmatrix}.
\]
The connecting map in the Tate exact sequence is $[v]\mapsto[av]$.
The off-diagonal relation
\[
 U'a+aU''=-S'b-bS''
\]
shows on each cycle that $\delta$ anticommutes with the connecting map
modulo boundaries.  All these cohomology groups are killed by two,
so it also commutes there.  Hence the even operator $u=\delta^2$ is
natural on the whole six-term periodic exact sequence.

Fix a polynomial $q\in\F_2[z]$ with $q(0)=1$.
On each finite-length Tate group take the nilpotent Fitting summand
for $q(u)$, namely $\ker(q(u)^N)$ for sufficiently large $N$.
This operation is exact and natural for maps commuting with $u$:
the Fitting decomposition into eventual kernel and image is a direct
sum, and both summands are preserved by every such map.  On the
chosen summand $u$ is invertible, by the coprimality of $z$ and $q$;
thus $\delta$ identifies its two parities.  The six-term exact
sequence becomes an exact three-cycle after this identification.
For any exact three-cycle of finite-length modules, the sum of the
three term lengths is twice the sum of the three image lengths.
Consequently
\[
 \nu_q(T)=\operatorname{length}_O
       \ker\bigl(q(\delta^2)^N:H^0\longrightarrow H^0\bigr)\pmod2
\]
is additive on the stated exact category.  It vanishes on every
scalar extension of an $O$-lattice, since $U=0$ there.

Choose distinct monic irreducibles $q_i\in\F_2[z]$ of odd degrees
$d_i$, with $q_i\ne z$; one may take one of each odd prime degree.
Let $D_i$ be the $0,1$ matrix lifting its companion matrix.  Define
\[
 T_i=\begin{pmatrix}
 I+\epsilon D_i&\pi I\\
 -h\epsilon D_i&-I-\epsilon D_i
 \end{pmatrix}.
\]
Direct multiplication gives $T_i^2=I$.  At $\epsilon=0$, its Tate
invariants are $H^0=k^{d_i}$ generated by the first block and
$H^1=k^{d_i}$ generated by the columns of $(I,-hI)$, where
$k=O/(\pi)$.  Indeed the norm image in the first block is
$(2,\pi)O^{d_i}=\pi O^{d_i}$, and the difference image in the other
eigenlattice is $\pi(I,-hI)O^{d_i}$.
The operator $\delta$ is $\overline D_i$ in both directions.  Hence
$u=\overline D_i^{\,2}$ and
\[
 \nu_{q_j}(T_i)=\begin{cases}1,&j=i,\\0,&j\ne i.\end{cases}
\]
Here squaring preserves the irreducible polynomial $q_i$ over
$\F_2$, and different $q_i$ stay coprime over $k$; no assumption that
$k$ is finite or perfect is needed.  The detected length is $d_i$,
which is odd.

Let $L_i$ be this lattice and put
$\gamma_i=[L_i]-d_i[A[C_2]]$.
Its reduction has a trivial free sublattice and sign free quotient,
so $\gamma_i$ lies in the stated reduction kernel.  The matrix
\[
 J_i=\begin{pmatrix}I&0\\-h(I+\epsilon D_i)&-I\end{pmatrix}
\]
is invertible and satisfies $J_iT_i=-T_iJ_i$.
The standard exact sequence
\[
 0\longrightarrow L_i\otimes\mathrm{sgn}
 \longrightarrow A[C_2]\otimes_A L_i
 \longrightarrow L_i\longrightarrow0
\]
therefore gives $2\gamma_i=0$.
The invariants $\nu_{q_j}$ prove independence and nonvanishing.
Only finitely many $q_j(\delta^2)$ can have nonzero nilpotent Fitting
summands on any given finite-length module, since the $q_j$ are
pairwise coprime.  Thus $(\nu_{q_j})_j$ takes values in the direct sum
and retracts the inclusion of the $\gamma_i$.
Finally every degree-two monomial in the trivial generators is a
scalar extension from $O$, whereas these classes are detected by
$\nu_{q_j}$.  This proves non-surjectivity of $\Phi_A$.
\end{proof}
